% historical-notes.tex — Historical context paragraph for insertion into the introduction.
% Intended length: 200–300 words.

The notion that a specific infrasonic frequency can induce involuntary
defecation---the so-called ``brown note''---has circulated in acoustic folklore
since at least the 1980s and entered mainstream culture through a 1999
\emph{South Park} episode and a 2005 \emph{MythBusters} segment that declared
the claim ``busted.''  The idea draws loosely on legitimate infrasound research,
but conflates two distinct physical phenomena: airborne acoustic loading and
mechanical whole-body vibration (WBV).

Early infrasound studies originated with Gavreau's experiments at CNRS
Levallois-Perret in the late 1950s, where a faulty ventilator at approximately
7\,Hz caused nausea among laboratory staff~\cite{Gavreau1966,Leventhall2007}.  Subsequent
NASA-funded work by Mohr et~al.\ exposed human subjects to infrasonic noise at
140--154\,dB\,SPL and reported chest-wall vibration, respiratory disruption, and
mild cardiovascular changes, but \emph{no gastrointestinal
effects}~\cite{Mohr1965}.  Von~Gierke's systematic review for
NASA~\cite{vonGierke1974} established exposure criteria that informed later ISO
standards.  Tandy and Lawrence~\cite{Tandy1998} demonstrated that a 19\,Hz
standing wave could induce visual disturbances and anxiety---real perceptual
effects, but far from the claimed bowel response.

Critical reviews by Leventhall~\cite{Leventhall2003,Leventhall2007} and
assessments of acoustic weapons by Altmann~\cite{Altmann1999} and Jauchem and
Cook~\cite{Jauchem2007} have concluded that infrasound is impractical as a
weapon and that claims of extreme physiological effects are exaggerated or
fabricated.  Meanwhile, the occupational health literature on WBV---codified in
ISO~2631~\cite{iso2631} and reviewed by Griffin~\cite{griffin1990handbook} and
Mansfield~\cite{mansfield2005human}---documents genuine gastrointestinal
disturbances among truck drivers, helicopter pilots, and construction equipment
operators exposed to mechanical vibration at 3--8\,Hz.  The present work takes
this well-attested mechanical coupling, rather than the unsupported acoustic
folklore, as its starting point.
